
\section{纤维中的伪全纯球面}\label{sec:8-4}

我们的目标是证明，截面\BookRecordIndex{T0246}模空间上的一个适当\BookRecordIndex{T0283}求值映射是\BookRecordIndex{T0337}伪循环。与
第~\ref{sec:6-7} 节讨论的图像情形一样，截面空间边界中的\BookRecordIndex{T0340}稳定映射由一个截面和若干位于
纤维中的球面组成。对一般\BookRecordIndex{T0279}纤维化 $\widetilde M\to\Sigma$，这些球面的分析与
$\widetilde M$ 是乘积的情形几乎完全相同，所以这里只陈述结果。

设
\[
  J\in\mathcal J^{\mathrm{Vert}}
  :=\mathcal J^{\mathrm{Vert}}(\widetilde M,\pi,\widetilde\omega)
\]
是\BookRecordIndex{T0306}竖直殆复结构，以 $\widetilde J$ 表示由 $J$ 与 $\widetilde\omega$ 确定的
$\widetilde M$ 上\BookRecordIndex{T0088}殆复结构，并设
$\widetilde A\in H_2(\widetilde M;\mathbb Z)$ 是满足
\[
  \pi_*\widetilde A=0
\]
的球面同调类。于是存在唯一的一族球面同调类
$A_z\in H_2(M_z;\mathbb Z)$（$z\in\Sigma$），使得
$\iota_{z*}A_z=\widetilde A$；这里
$\iota_z:M_z\to\widetilde M$ 表示纤维的包含映射。条件
$\pi_*\widetilde A=0$ 蕴含：对每个 $\widetilde J$-全纯球面
$\widetilde v:S^2\to\widetilde M$，复合 $\pi\circ\widetilde v$ 都是常值映射。此外，
\[
  c_1^{\mathrm{Vert}}(\widetilde A)
  =\langle c_1(TM_z),A_z\rangle
  =\langle c_1(T\widetilde M),\widetilde A\rangle
  =c_1(\widetilde A).
\]

以下考察\BookRecordIndex{T0246}模空间
\[
  \mathcal M^*(\widetilde A;\widetilde J)
  :=\bigl\{(z,v)\mid z\in\Sigma,\ v\in\mathcal M^*(A_z;J_z)\bigr\},
\]
它由各纤维中表示类 $\widetilde A$ 的单纯 $\widetilde J$-全纯球面组成。与
$\widetilde J$-全纯球面 $\widetilde v:S^2\to\widetilde M$ 相联系的是
第~\ref{sec:3-1} 节引入的\BookRecordIndex{T0210}柯西--黎曼算子
\[
  D_{\widetilde v}=D_{\widetilde J,\widetilde v}:
  \Omega^0(S^2,\widetilde v^*T\widetilde M)
  \longrightarrow
  \Omega^{0,1}(S^2,\widetilde v^*T\widetilde M).
\]
需要强调的是，\BookRecordIndex{T0088}殆复结构 $\widetilde J$ 和算子 $D_{\widetilde v}$ 都依赖于\BookRecordIndex{T0224}联络形式
$\widetilde\omega$，尽管\BookRecordIndex{T0246}模空间
$\mathcal M^*(\widetilde A;\widetilde J)$ 只依赖于\BookRecordIndex{T0306}竖直殆复结构 $J$。不过，正如
注~\ref{rem:6-7-8}，对每个（竖直的）$\widetilde J$-全纯球面
$\widetilde v:S^2\to\widetilde M$，$D_{\widetilde v}$ 的核与这个选择无关。具体地说，
每个这样的球面都满足 $\pi\circ\widetilde v\equiv\mathrm{const}$，因而可考察限制算子
\[
  D_{\widetilde v}^{\mathrm{Vert}}:
  \bigl\{\widetilde\xi\in\Omega^0(S^2,\widetilde v^*T\widetilde M)\mid
  d\pi(\widetilde v)\widetilde\xi\equiv\mathrm{const}\bigr\}
  \longrightarrow
  \Omega^{0,1}(S^2,\widetilde v^*T^{\mathrm{Vert}}\widetilde M).
\]
公式
\[
  d\pi(\widetilde v)\circ D_{\widetilde v}
  =\bar\partial\circ d\pi(\widetilde v)
\]
表明这个算子良定，它与 $D_{\widetilde v}$ 有相同的核，并且两个算子的余核同构。因此
$D_{\widetilde v}$ 与 $D_{\widetilde v}^{\mathrm{Vert}}$ 是指标相同的\BookRecordIndex{T0128}弗雷德霍姆算子，且
\[
  \operatorname{index}D_{\widetilde v}^{\mathrm{Vert}}
  =\operatorname{index}D_{\widetilde v}
  =2n+2+2c_1(\widetilde A).
\]
而且，$D_{\widetilde v}^{\mathrm{Vert}}$ 满射当且仅当 $D_{\widetilde v}$ 满射。留给读者
验证 $D_{\widetilde v}^{\mathrm{Vert}}$ 与 $\widetilde\omega$ 的选择无关。注意，在平凡丛
情形，$D_{\widetilde v}^{\mathrm{Vert}}$ 正是式~\eqref{eq:6-7-8} 的算子 $D_{z,v}$，以上
断言已经在注~\ref{rem:6-7-8} 中证明。

设 $\widetilde A\in H_2(\widetilde M;\mathbb Z)$ 满足
$\pi_*\widetilde A=0$。若对每个表示类 $\widetilde A$ 的单纯 $\widetilde J$-全纯球面
$\widetilde v:S^2\to\widetilde M$，算子 $D_{\widetilde v}$ 都满射，就称\BookRecordIndex{T0306}竖直殆复结构
$J\in\mathcal J^{\mathrm{Vert}}$ \emph{关于 $\widetilde A$ 正则}。关于
$\widetilde A$ 正则的\BookRecordIndex{T0306}竖直殆复结构之集记为
\[
  \mathcal J_{\mathrm{reg}}^{\mathrm{Vert}}(\widetilde A).
\]
下述定理的证明几乎与定理~\ref{thm:3-1-6} 的证明逐字相同，故从略。

\begin{theorem}\label{thm:8-4-1}
设 $\widetilde A\in H_2(\widetilde M;\mathbb Z)$ 满足
$\pi_*\widetilde A=0$。
\begin{enumerate}[label=\textup{(\roman*)},leftmargin=3em]
\item 若 $J\in\mathcal J_{\mathrm{reg}}^{\mathrm{Vert}}(\widetilde A)$，则
$\mathcal M^*(\widetilde A;\widetilde J)$ 是光滑定向流形，维数为
\[
  \dim\mathcal M^*(\widetilde A;\widetilde J)
  =2n+2+2c_1(\widetilde A).
\]
\item 集合 $\mathcal J_{\mathrm{reg}}^{\mathrm{Vert}}(\widetilde A)$ 是
$\mathcal J^{\mathrm{Vert}}$ 中的\BookRecordIndex{T0300}剩余集。
\end{enumerate}
\end{theorem}

现在取
$J_0,J_1\in\mathcal J_{\mathrm{reg}}^{\mathrm{Vert}}(\widetilde A)$，并设
\[
  [0,1]\longrightarrow\mathcal J^{\mathrm{Vert}},
  \qquad\lambda\longmapsto J_\lambda
\]
是从 $J_0$ 到 $J_1$ 的光滑同伦。对任意这样的同伦考察\BookRecordIndex{T0246}模空间
\[
  \mathcal W^*(\widetilde A;\{\widetilde J_\lambda\}_\lambda)
  :=\bigl\{(\lambda,\widetilde v)\mid0\le\lambda\le1,\
  \widetilde v\in\mathcal M^*(\widetilde A;\widetilde J_\lambda)\bigr\}.
\]
若 $J_0,J_1\in\mathcal J_{\mathrm{reg}}^{\mathrm{Vert}}(\widetilde A)$，且对每个
$\lambda\in[0,1]$ 和每个表示类 $\widetilde A$ 的单纯
$\widetilde J_\lambda$-全纯球面 $\widetilde v:S^2\to\widetilde M$ 都有
\[
  \Omega^{0,1}(S^2,\widetilde v^*T^{\mathrm{Vert}}\widetilde M)
  =\operatorname{im}D^{\mathrm{Vert}}_{J_\lambda,\widetilde v}
   +\mathbb R(\partial_\lambda J_\lambda)d\widetilde v\circ i,
\]
就称同伦 $\{J_\lambda\}_\lambda$ \emph{正则}。\BookRecordIndex{T0407}正则同伦之空间记为
$\mathcal J_{\mathrm{reg}}^{\mathrm{Vert}}(\widetilde A;J_0,J_1)$，所有光滑同伦之空间记为
$\mathcal J^{\mathrm{Vert}}(J_0,J_1)$。下面的定理几乎与定理~\ref{thm:3-1-8} 逐字相同，
证明可从略。

\begin{theorem}\label{thm:8-4-2}
设 $J_0,J_1\in\mathcal J_{\mathrm{reg}}^{\mathrm{Vert}}(\widetilde A)$，且
$\widetilde A\in H_2(\widetilde M;\mathbb Z)$ 满足 $\pi_*\widetilde A=0$。
\begin{enumerate}[label=\textup{(\roman*)},leftmargin=3em]
\item 若
$\{J_\lambda\}_\lambda\in
\mathcal J_{\mathrm{reg}}^{\mathrm{Vert}}(\widetilde A;J_0,J_1)$，则
$\mathcal W^*(\widetilde A;\{\widetilde J_\lambda\}_\lambda)$ 是带边光滑定向流形，且
\[
  \partial\mathcal W^*(\widetilde A;\{\widetilde J_\lambda\}_\lambda)
  =\mathcal M^*(\widetilde A;\widetilde J_1)
   -\mathcal M^*(\widetilde A;\widetilde J_0).
\]
\item 集合
$\mathcal J_{\mathrm{reg}}^{\mathrm{Vert}}(\widetilde A;J_0,J_1)$ 是
$\mathcal J^{\mathrm{Vert}}(J_0,J_1)$ 中的\BookRecordIndex{T0300}剩余集。
\end{enumerate}
\end{theorem}

\begin{example}\label{ex:8-4-3}
令 $M=S^2\times S^2$，赋予乘积辛结构
\[
  \omega_\lambda:=\lambda\pi_1^*\sigma+\pi_2^*\sigma,
  \qquad\lambda>1,
\]
其中 $\sigma$ 是 $S^2$ 上的标准面积形式，$\pi_i$ 是到第 $i$ 个因子的投影。考察\BookRecordIndex{T0279}纤维化
$\pi_1:M\to S^2$。$M$ 上有一个圆周作用
$\{\psi_t\}_{t\in S^1}$，它旋转每个纤维，并固定对角线
\[
  \{(x,x)\mid x\in S^2\subset\mathbb R^3\}
\]
和反对角线
\[
  \{(x,-x)\mid x\in S^2\subset\mathbb R^3\}
\]
上的点。（这个作用的公式见下文式~\eqref{eq:9-7-1}。）我们将在命题~\ref{prop:9-7-2}\textup{(ii)} 中证明，
当 $\lambda>1$ 时，环路 $\{\psi_t\}$ 在
$\Ham(M,\omega_\lambda)$ 中同痕于环路 $\{\psi_t^\lambda\}$。按照
注~\ref{rem:8-2-11} 前的讨论，用此环路构造纤维丛
\sourcefootnote{原书把纤维写成 $(M,\omega)$；但本例只定义了 $\omega_\lambda$，且构造
所用环路明确位于 $\Ham(M,\omega_\lambda)$，故纤维的\BookRecordIndex{T0363}辛形式应为 $\omega_\lambda$。}
\[
  (M,\omega_\lambda)\longrightarrow\widetilde M\longrightarrow S^2.
\]
反对角线在 $M$ 中的类 $\widetilde A$ 在 $M$ 中的自交数为 $-2$，Chern 数为 $0$，故对
通有的被 $\omega_\lambda$ \BookRecordIndex{T0372}驯服的 $J$，它没有 $J$-全纯代表元（比较
例~\ref{ex:7-2-2}）。不过，Le--Ono \srcite{236} 证明了
\[
  c_1(\widetilde M)(\widetilde A)=1,
  \qquad
  \operatorname{GW}^{\widetilde M}_{\widetilde A,0}=1.
\]
他们用参数\BookRecordIndex{T0143}格罗莫夫--威滕不变量解释了这个结果；另见 Buse \srcite{54}。这里的解析要点
是：若曲线
\[
  \widetilde v=\iota\circ v:
  S^2\xrightarrow{v}M\xhookrightarrow{\iota}\widetilde M
\]
位于单个纤维中，则上面描述的\BookRecordIndex{T0346}线性化算子
$D_{\widetilde v}^{\mathrm{Vert}}$ 的像，与 $M$ 中相应算子 $D_v$ 的像有所不同；差别来自
族 $J_z$ 关于基变量 $z\in\Sigma$ 的导数项。精确公式见式~\eqref{eq:6-7-9}。
\end{example}

\begin{remark}\label{rem:8-4-4}
我们正在发展的方法稍后将用于构造 \BookRecordIndex{T0047}Seidel 表示（见例 8.6.8）；它的一个应用是证明，
\BookRecordIndex{T0156}哈密顿纤维化 $(\widetilde M,\widetilde\omega)\to\Sigma$ 的总空间之有理同调是乘积：
\[
  H_*(\widetilde M;\mathbb Q)
  \cong H_*(\Sigma;\mathbb Q)\otimes H_*(M;\mathbb Q).
\]
当 $(M,\omega)$ \BookRecordIndex{T0058}半正时，这一点由命题 11.4.5 证明，其思路来自
Lalonde--McDuff--Polterovich \srcite{228}。一般情形的证明见 McDuff \srcite{269}。
因此，有理类 $A_z\in H_2(M_z;\mathbb Q)$ 由 $\widetilde A$ 唯一确定，并可等同于单个类
$A\in H_2(M;\mathbb Q)$。（这里用到了如下事实：$\widetilde M\to\Sigma$ 的每个纤维
$M_z$ 都能由一个辛微分同胚与参考纤维等同，而这个辛微分同胚在\BookRecordIndex{T0158}哈密顿同痕的意义下良定。）

上述关于 $H_*(\widetilde M;\mathbb Q)$ 结构的断言对曲面上的一般\BookRecordIndex{T0200}局部哈密顿纤维化并不
成立。例如，考虑第~\ref{sec:8-2} 节中的\BookRecordIndex{T0381}映射圆 $Y_f\to\mathbb T^2$，其中构造它的映射
$f:S^2\times S^2\to S^2\times S^2$ 交换两个球面。

本节的思想可以推广到更广的框架：若只关心纤维中的球面，就不必把\BookRecordIndex{T0279}纤维化限制在曲面上，
事实上基底甚至不必是\BookRecordIndex{T0358}辛流形。若干很有意思的应用见 Seidel \srcite{372}，特别是其中的
例 1.13。
\end{remark}
